% Shared body for the CS 138 proposal and final report.
%
% Do not compile this file directly. Build proposal.tex or final.tex instead
% (or run `make proposal` / `make final`). Those files set \ifproposal and then
% \input this one, so both documents come from a single source.
%
% Sections wrapped in \ifproposal ... \else ... \fi differ between the two
% documents; everything else appears in both.

\usepackage[utf8]{inputenc}
\usepackage[numbers]{natbib}
\usepackage{graphicx}
\usepackage{amsmath,amssymb}
\usepackage{hyperref}

\title{A Very Cool and Catchy Title}
\author{Jake \and Harris Davidson \and Rachel Ellinger}
\date{\vspace{-1em}}

\begin{document}

\maketitle

\ifproposal\else
\begin{abstract}
A one paragraph high level description of the work. The abstract is typically
the first thing someone would read from the paper before deciding whether to
continue reading and hence, serves as an advertisement to the reader to read
the whole paper.
\end{abstract}
\fi

\section{Introduction}

% DRAFT. Requirements (from the course template): describe the problem in a
% non-technical manner (no math or equations) and motivate it, i.e., explain
% why it is important.

When you hear a sound you cannot place, you rarely stay still. You turn your
head, listen again, and turn some more until the source snaps into place.
Listening is an active process: where you point your ears determines what you
hear, and what you hear determines where you should point them next. This
project asks whether an artificial agent can learn that same habit, and learn
to do it well.

We study a simple version of the problem. Two microphones sit on a motorized
``head'' that can turn left and right, and somewhere in the room is a sound
source, which may be moving. The agent's job is to work out where the sound is
coming from as accurately and as quickly as possible. Holding still is often
not enough: with only two microphones, a sound from the front and a sound from
directly behind can produce nearly the same signals, so a single reading
cannot tell them apart. Turning the head and listening again breaks the tie,
but every turn costs time, and a moving source will not wait while the agent
deliberates. The agent therefore faces a planning problem: deciding how to turn
so that it collects the information it needs to locate the sound, without
wasting effort on movements that tell it little, and knowing when it has heard
enough to commit to an answer.

This makes the task a natural fit for reinforcement learning. Rather than
hand-coding a search strategy, we let the agent learn one through trial and
error, rewarding it for naming the right direction and charging it for every
turn it takes and every wrong guess. We then compare what it learns against
simple strategies, such as turning at random or following a fixed routine that
turns until the sound is centered and then checks front against back.

We train the agent in a simulated room, which lets us control exactly what it
hears and run far more trials than would be practical with real equipment. We
start with the simplest version of the problem, a quiet room with no echoes and
a single stationary source playing short bursts of static, and then add realism
one step at a time: echoes, competing sounds, small imperfections in the
microphones and motor, a moving source, and a second direction of rotation. We
are also building a small set of physical heads, each a pair of microphones on
a motor, to test whether behavior learned in simulation holds up in a real
room and to demonstrate the result.

The problem matters well beyond this setup. Robots that work alongside people
need to turn toward someone who speaks to them; search-and-rescue systems may
need to find a person by the sound they make; and devices such as smart
speakers, security cameras, and hearing aids all benefit from knowing where a
sound is coming from. In each case, sensors are limited and time matters, so
the ability to choose where to listen next is as important as the ability to
interpret what is heard. As stretch goals, we also plan to examine whether a
learned strategy transfers to new conditions, and to extend the task to
multiple agents in a setting modeled on the game Marco Polo.

\section{Background and Related Work}

Following, you should provide the necessary background and discuss related
work in the RL literature. This section should also be about a page.
Citations should be in BibTeX format \citep{thrun2005probabilistic}.

\section{Technical Approach / Methodology / Theoretical Framework}

\ifproposal
Describe how you will approach the problem and its technical formulation.
Feel free to re-state the basic RL formulas (e.g., if using Q-learning, state
the update rule or the formula for what the Q function approximates).
\else
A detailed description of your problem (with math, notation, algorithms,
figures, etc.). Use footnotes to cite links to your code or
videos\footnote{All developed source code for this project is available at
\url{https://github.com/djake/sound-localization-rl}.}

\subsection{Notation and Problem Formulation}

How did you formalize the problem using mathematical notation?

\subsection{Tasks}

Subsections are useful for breaking down the problem into sub-parts. For
example, you could break down the tasks for your project and list them one by
one.
\fi

\ifproposal
\section{Experiments and Evaluation}

Describe how you will evaluate your approach/solution. What constitutes
success? What metrics will you use? Do you have any preliminary hypothesis
that you plan to test? Also, describe the RL domain or environment you plan to
use.

\section{Timeline and Individual Responsibilities}

State the timeline in terms of weeks and milestones you want to achieve. If
working on a team, state what the individual responsibilities are at this
point (i.e., who is going to do what, these may of course change over the
course of the project) \cite{short2010no}.
\else
\section{Experimental Results / Technical Demonstration}

A description of how you evaluated or demonstrated your
solution.\footnote{A video of the agent doing x y z is available at \ldots}

\section{Conclusion and Future Work}

A high level summary of what was accomplished, along with a discussion on
limitations and avenues for future work (typically 2 to 3 paragraphs).
\fi

\bibliographystyle{plain}
\bibliography{references}

\end{document}
